\documentclass[letterpaper]{article}
\usepackage[margin=1in]{geometry}
\usepackage{times}
\usepackage[T1]{fontenc}
\usepackage[utf8]{inputenc}
\usepackage{amsmath,amssymb,amsthm}
\usepackage{booktabs}
\usepackage{array}
\usepackage{longtable}
\usepackage{xcolor}
\usepackage[colorlinks=true,linkcolor=blue,citecolor=green!60!black,urlcolor=blue]{hyperref}
\usepackage{microtype}
\usepackage{enumitem}
\usepackage{float}
\usepackage{caption}

\newtheorem{theoreme}{Th\'eor\`eme}
\newtheorem{corollaire}{Corollaire}[theoreme]
\newtheorem{lemme}{Lemme}
\newtheorem{definition}{D\'efinition}

\title{%
\textbf{Inf\'erence d'\'Etat de R\`egles (RSI) ---}\\
\textbf{Mat\'eriau Suppl\'ementaire}\\[0.4em]
\large\textit{Annexes A, B, C}
}
\author{Abdou-Raouf Atarmla}
\date{2026}

\begin{document}
\maketitle
\thispagestyle{empty}

\noindent Ce document contient le mat\'eriau suppl\'ementaire associ\'e
au papier \emph{Inf\'erence d'\'Etat de R\`egles (RSI) : Un Cadre
Bay\'esien pour la Surveillance de la Conformit\'e dans les Domaines
R\'egis par des R\`egles}. Il comprend le dictionnaire du dataset
(Annexe~A), les preuves compl\`etes des th\'eor\`emes (Annexe~B), et
les d\'etails de reproductibilit\'e (Annexe~C).

\tableofcontents
\newpage

%% ─────────────────────────────────────────────────────────────────
\appendix

\section{Dictionnaire du Dataset : RSI-Togo-Fiscal-Synthetic v1.0}
\label{app:dataset}

\subsection{Vue d'ensemble}

Le dataset contient 2~000 entreprises synth\'etiques couvrant deux
p\'eriodes r\'eglementaires (2022--2024 et 2025), g\'en\'er\'ees \`a
partir des r\`egles fiscales r\'eelles de l'OTR. Chaque ligne est
structur\'ee en quatre couches~:

\begin{enumerate}[leftmargin=2em]
  \item \textbf{Caract\'eristiques des entreprises}~: attributs
        observables de l'entit\'e
  \item \textbf{V\'erit\'e terrain latente} (\texttt{gt\_*})~: le
        vrai \'etat de la r\`egle $(a_i, c_i, \delta_i)$, jamais
        accessible \`a RSI durant l'inf\'erence
  \item \textbf{Observations bruit\'ees} (\texttt{obs\_*})~: ce
        qu'un syst\`eme fiscal observe r\'eellement
  \item \textbf{\'Etiquettes binaires} (\texttt{label\_*})~: labels
        de conformit\'e d\'eriv\'es pour l'\'evaluation des baselines
\end{enumerate}

\subsection{R\'ef\'erence des colonnes}

\begin{longtable}{lp{2.2cm}p{7.5cm}}
\toprule
\textbf{Colonne} & \textbf{Type} & \textbf{Description} \\
\midrule
\endfirsthead
\toprule
\textbf{Colonne} & \textbf{Type} & \textbf{Description} \\
\midrule
\endhead
\midrule
\multicolumn{3}{r}{\small Suite page suivante} \\
\endfoot
\bottomrule
\endlastfoot

\multicolumn{3}{l}{\textit{Identifiants}} \\[2pt]
\texttt{enterprise\_id} & cha\^ine & Identifiant unique de l'entreprise (ex. ENT\_00001) \\
\texttt{period} & cha\^ine & P\'eriode r\'eglementaire~: \texttt{2022\_2024} ou \texttt{2025} \\[6pt]

\multicolumn{3}{l}{\textit{Caract\'eristiques des entreprises}} \\[2pt]
\texttt{sector} & cha\^ine & Secteur \'economique (commerce, services, BTP, industrie, agriculture, restauration, transport) \\
\texttt{region} & cha\^ine & R\'egion g\'eographique au Togo (Lom\'e, Sokod\'e, Kara, etc.) \\
\texttt{ca\_segment} & cha\^ine & Segment de CA~: \texttt{informal} ($<$30M), \texttt{rsi} (30--100M), \texttt{reel} ($\geq$100M) \\
\texttt{n\_employes} & int & Nombre r\'eel d'employ\'es (distribution log-normale) \\[6pt]

\multicolumn{3}{l}{\textit{V\'erit\'e terrain latente (gt\_*) --- inaccessible \`a RSI durant l'inf\'erence}} \\[2pt]
\texttt{gt\_ca\_reel} & flottant & Chiffre d'affaires annuel r\'eel en FCFA \\
\texttt{gt\_benefice\_reel} & flottant & B\'en\'efice annuel r\'eel en FCFA \\
\texttt{gt\_tva\_active} & int & 1 si la r\`egle TVA s'applique (CA $\geq$ seuil), 0 sinon \\
\texttt{gt\_tva\_compliance} & flottant & Taux de conformit\'e TVA r\'eel $c_i \in [0,1]$, tir\'e d'un prior Beta \\
\texttt{gt\_tva\_due} & flottant & TVA th\'eorique due (FCFA) \\
\texttt{gt\_is\_active} & int & 1 si la r\`egle IS s'applique (CA $\geq$ 100M FCFA) \\
\texttt{gt\_is\_compliance} & flottant & Taux de conformit\'e IS r\'eel $c_i \in [0,1]$ \\
\texttt{gt\_is\_due} & flottant & IS effectif d\^u (max de l'IS et de l'IMF) en FCFA \\
\texttt{gt\_tpu\_active} & int & 1 si la r\`egle TPU s'applique (CA $<$ 30M FCFA) \\
\texttt{gt\_tpu\_compliance} & flottant & Taux de conformit\'e TPU r\'eel $c_i \in [0,1]$ \\
\texttt{gt\_irpp\_susp\_active} & int & 1 si la suspension IRPP s'applique (salaire $<$ 900k, p\'eriode 2023) \\[6pt]

\multicolumn{3}{l}{\textit{Observations bruit\'ees (obs\_*) --- entr\'ee de RSI}} \\[2pt]
\texttt{obs\_ca\_declare} & flottant & CA d\'eclar\'e. G\'en\'er\'e comme $\hat{x} = x \cdot \beta \cdot \varepsilon$ o\`u $\beta \sim \text{Beta}(7,3)$ capture la sous-d\'eclaration syst\'ematique et $\varepsilon \sim \mathcal{N}(1, 0{,}045)$ le bruit de mesure. Ratio moyen $\mathbb{E}[\hat{x}/x] \approx 0{,}70$. \\
\texttt{obs\_tva\_declaree} & flottant & Montant de TVA d\'eclar\'e. Z\'ero si l'entit\'e ne d\'eclare pas~; distribu\'e log-normalement autour de la TVA th\'eorique sinon. \texttt{NaN} avec probabilit\'e 0,18 (donn\'ee manquante). \\
\texttt{obs\_tva\_assujetti\_declare} & bool & Si l'entit\'e s'est auto-d\'eclar\'ee assujettie \`a la TVA \\
\texttt{obs\_is\_declare} & flottant & IS d\'eclar\'e. \texttt{NaN} avec probabilit\'e 0,18. \\
\texttt{obs\_benefice\_declare} & flottant & B\'en\'efice d\'eclar\'e. Typiquement 70\% du b\'en\'efice r\'eel avec bruit. \texttt{NaN} avec probabilit\'e 0,18. \\
\texttt{obs\_retard\_paiement\_jours} & int & D\'elai de paiement en jours. G\'en\'er\'e depuis $\text{Exponentielle}(\lambda)$ o\`u $\lambda = 5 / (c_{\text{r\'egime}} + 0{,}1)$, avec $c_{\text{r\'egime}}$ la conformit\'e du r\'egime applicable. Non-conformit\'e \'elev\'ee $\Rightarrow$ d\'elais plus longs. \\
\texttt{obs\_has\_compte\_bancaire} & bool & Si l'entit\'e poss\`ede un compte bancaire formel. $P(\text{compte}) = 0{,}75 \cdot c + 0{,}20 \cdot (1-c)$ o\`u $c$ est la conformit\'e du r\'egime. \\
\texttt{obs\_utilise\_facturation\_electronique} & bool & Si l'entit\'e utilise la facturation \'electronique \\
\texttt{obs\_a\_ete\_audite} & bool & Si l'entit\'e a \'et\'e audit\'ee sur la p\'eriode (taux de base 8\%) \\
\texttt{obs\_n\_employes\_declare} & int & Nombre d'employ\'es d\'eclar\'e ($\approx$~85\% de la vraie valeur) \\
\texttt{obs\_ratio\_sous\_declaration} & flottant & Ratio CA d\'eclar\'e / CA r\'eel ($\hat{x}/x$) \\
\texttt{obs\_tva\_missing} & bool & Vrai si \texttt{obs\_tva\_declaree} est manquant \\
\texttt{obs\_is\_missing} & bool & Vrai si \texttt{obs\_is\_declare} est manquant \\
\texttt{obs\_benefice\_missing} & bool & Vrai si \texttt{obs\_benefice\_declare} est manquant \\[6pt]

\multicolumn{3}{l}{\textit{\'Etiquettes binaires (label\_*) --- utilis\'ees uniquement pour l'\'evaluation des baselines}} \\[2pt]
\texttt{label\_tva\_non\_conforme} & int & 1 si TVA active et $c_{\text{TVA}} < 0{,}5$ \\
\texttt{label\_is\_non\_conforme} & int & 1 si IS actif et $c_{\text{IS}} < 0{,}5$ \\
\texttt{label\_any\_non\_conforme} & int & 1 si au moins l'un de~: TVA non conforme, IS non conforme, ou conformit\'e TPU $< 0{,}3$ \\

\end{longtable}

\subsection{Mod\`ele de sous-d\'eclaration}

Le ratio de sous-d\'eclaration est le m\'ecanisme de bruit principal,
refl\'etant le comportement fiscal r\'eel en Afrique francophone.
Pour chaque entreprise~:

\[
  \hat{x} = x \cdot \underbrace{\beta}_{\text{syst\'ematique}} \cdot
  \underbrace{\varepsilon}_{\text{bruit}}, \quad
  \beta \sim \text{Beta}(7, 3), \quad
  \varepsilon \sim \mathcal{N}(1,\, 0{,}045^2)
\]

La distribution Beta(7,3) a une moyenne $7/10 = 0{,}70$ et concentre
sa masse entre 0,55 et 0,85, conform\'ement aux estimations empiriques
des \'etudes fiscales en Afrique subsaharienne. Le terme de bruit
gaussien capture les approximations comptables et les arrondis.

\subsection{\'Ev\'enement de changement r\'eglementaire}

La p\'eriode \texttt{2025} encode le rel\`evement du seuil TVA de
60M \`a 100M FCFA par la Loi n°2024-007 (30 d\'ecembre 2024).
Les entreprises avec un CA r\'eel dans $[60\text{M}, 100\text{M})$
FCFA passent de \texttt{gt\_tva\_active=1} (p\'eriode 2022--2024)
\`a \texttt{gt\_tva\_active=0} (p\'eriode 2025), fournissant une
exp\'erience naturelle propre pour valider le Th\'eor\`eme~1
(adaptabilit\'e $O(1)$).

\subsection{Structure en loi de puissance}

Le CA est tir\'e d'une distribution log-normale par segment, dont
l'agr\'egat approxime une loi de puissance---la signature universelle
des distributions de taille des firmes dans les \'economies de
march\'e~\cite{gabaix2009power}. La structure en trois segments
reproduit la concentration typique des march\'es africains~:

\begin{center}
\begin{tabular}{lcc}
\toprule
\textbf{Segment} & \textbf{Plage de CA} & \textbf{Part} \\
\midrule
Informel (TPU) & $<$ 30M FCFA & 60\% \\
Interm\'ediaire (RSI) & 30M -- 100M FCFA & 25\% \\
Grand contribuable (IS/TVA) & $\geq$ 100M FCFA & 15\% \\
\bottomrule
\end{tabular}
\end{center}

%% ─────────────────────────────────────────────────────────────────
\section{Preuves Compl\`etes des Th\'eor\`emes}
\label{app:preuves}

\subsection{Preuve du Th\'eor\`eme 1 : Adaptabilit\'e $O(1)$}

\begin{theoreme}[Adaptabilit\'e $O(1)$ --- complet]
Pour toute mise \`a jour r\'eglementaire $\mathcal{U}_k : \Theta_k
\to \Theta_k'$, le co\^ut de mise \`a jour RSI est
$\mathcal{C}^{\text{RSI}}(\mathcal{U}_k) = O(1)$, ind\'ependamment
de $|\mathbf{D}|$ et de $n$. De plus, le posterior mis \`a jour est
correctement normalis\'e, sans risque d'explosion de variance.
\end{theoreme}

\begin{proof}
\textbf{\'Etape 1 : Factorisation.}
L'a priori RSI se factorise comme~:
\[
  P(\mathbf{S}) = \prod_{i=1}^n P(s_i \mid \Theta_i)
\]
Apr\`es la mise \`a jour $\mathcal{U}_k$, le nouvel a priori est~:
\[
  P'(\mathbf{S}) = P(s_k' \mid \Theta_k') \cdot \prod_{i \neq k} P(s_i \mid \Theta_i)
\]

\textbf{\'Etape 2 : Correction du posterior.}
Par le th\'eor\`eme de Bayes, le posterior avant mise \`a jour est~:
\[
  P(\mathbf{S} \mid \mathbf{D}) = \frac{P(\mathbf{D} \mid \mathbf{S})
  \cdot P(\mathbf{S})}{P(\mathbf{D})}
\]
Apr\`es la mise \`a jour~:
\begin{align*}
  P'(\mathbf{S} \mid \mathbf{D})
  &= \frac{P(\mathbf{D} \mid \mathbf{S}) \cdot P'(\mathbf{S})}{P'(\mathbf{D})} \\
  &\propto P(\mathbf{S} \mid \mathbf{D}) \cdot
    \underbrace{\frac{P(s_k' \mid \Theta_k')}{P(s_k \mid \Theta_k)}}_{\rho_k(\Theta_k, \Theta_k')}
\end{align*}

\textbf{\'Etape 3 : Complexit\'e.}
Le facteur de correction $\rho_k$ ne d\'epend que des formes
param\'etriques de $P(s_k \mid \Theta_k)$ et $P(s_k' \mid \Theta_k')$.
L'\'evaluation de deux distributions param\'etriques (Bernoulli,
Beta, Gaussienne) requiert un nombre constant d'op\'erations
arithm\'etiques, d'o\`u $O(1)$.

\textbf{\'Etape 4 : Stabilit\'e num\'erique.}
Le ratio $\rho_k$ est appliqu\'e au \emph{noyau} (forme non
normalis\'ee) du posterior, et non \`a sa valeur absolue.
La constante de normalisation est mise \`a jour comme~:
\[
  P'(\mathbf{D}) = P(\mathbf{D}) \cdot
  \mathbb{E}_{P(\mathbf{S} \mid \mathbf{D})}[\rho_k]
\]
Sous les familles conjugu\'ees utilis\'ees dans RSI (Beta-Binomiale
pour $c_i$, Bernoulli pour $a_i$, Gaussienne pour $\delta_i$), cette
esp\'erance est disponible en forme ferm\'ee. Le posterior reste donc
une distribution de probabilit\'e proprement normalis\'ee apr\`es la
mise \`a jour, sans explosion de variance.
\end{proof}

\begin{corollaire}[Co\^ut cumulatif sur $T$ mises \`a jour]
Sur $T$ mises \`a jour r\'eglementaires, les co\^uts totaux scalent
comme~:
\begin{center}
\begin{tabular}{lc}
\toprule
\textbf{M\'ethode} & \textbf{Co\^ut total} \\
\midrule
RSI & $O(T)$ \\
ML supervis\'e (r\'e-entra\^inement complet) & $O(T \cdot |\mathbf{D}| \cdot E)$ \\
PSL / MLN & $O(T \cdot |\mathbf{D}| \cdot n)$ \\
\bottomrule
\end{tabular}
\end{center}
o\`u $E$ d\'enote le nombre d'\'epoques d'entra\^inement. RSI domine
dans les environnements \`a forte fr\'equence r\'eglementaire.
\end{corollaire}

\subsection{Preuve du Th\'eor\`eme 2 : Consistance de Bernstein-von Mises}

\textbf{V\'erification des conditions de r\'egularit\'e.}

\begin{itemize}[leftmargin=2em]
  \item[\textbf{C1}] \textit{Identifiabilit\'e.}
        $P(\mathbf{D} \mid \mathbf{S}) \neq P(\mathbf{D} \mid \mathbf{S}')$
        pour $\mathbf{S} \neq \mathbf{S}'$.
        \textit{V\'erification~:} Chaque composante $(a_i, c_i, \delta_i)$
        influence la vraisemblance via des m\'ecanismes distincts.
        Des \'etats distincts produisent des fonctions de vraisemblance
        distinctes.

  \item[\textbf{C2}] \textit{A priori positif en la v\'erit\'e.}
        $P(\mathbf{S}^*) > 0$.
        \textit{V\'erification~:} L'a priori est un produit de Bernoulli,
        Beta et Gaussiennes, toutes \`a densit\'e strictement positive sur
        leurs supports. Donc $P(\mathbf{S}) > 0$ pour tout $\mathbf{S}
        \in \mathcal{S}$.

  \item[\textbf{C3}] \textit{R\'egularit\'e de la vraisemblance.}
        $\log P(\mathbf{D} \mid \mathbf{S})$ est deux fois
        diff\'erentiable dans les composantes continues $(c_i, \delta_i)$.
        \textit{V\'erification~:} Les fonctions de vraisemblance
        (log-gaussienne, exponentielle, Beta-Binomiale) sont lisses en
        leurs param\`etres. La diff\'erentiabilit\'e est h\'erit\'ee
        par le produit.

  \item[\textbf{C4}] \textit{Information de Fisher finie.}
        $\mathcal{I}(\mathbf{S}^*)$ est finie et d\'efinie positive.
        \textit{V\'erification~:} Sous C1 et C3, et sous r\'eserve
        qu'aucune r\`egle ne soit d\'eg\'en\'er\'ee (i.e., $c_i \notin
        \{0,1\}$ et $a_i \in (0,1)$), la matrice d'information de Fisher
        est d\'efinie positive.
\end{itemize}

\begin{theoreme}[Consistance BvM --- complet]
Sous les conditions C1--C4, quand $m \to \infty$~:
\[
  \sqrt{m}\bigl(\hat{\mathbf{S}}_m - \mathbf{S}^*\bigr)
  \xrightarrow{d} \mathcal{N}\!\left(0,\, \mathcal{I}(\mathbf{S}^*)^{-1}\right)
\]
\end{theoreme}

\begin{proof}
\textbf{\'Etape 1 : Convergence par la LGN.}
Par la loi des grands nombres, pour des observations i.i.d.~:
\[
  \frac{1}{m} \log P(\mathbf{D}_m \mid \mathbf{S})
  \xrightarrow{p.s.} \mathbb{E}_{\mathbf{S}^*}[\log P(d \mid \mathbf{S})]
  = -\mathrm{KL}(P(\cdot \mid \mathbf{S}^*) \| P(\cdot \mid \mathbf{S}))
  + \text{cst}
\]
Ceci est maximis\'e en $\mathbf{S} = \mathbf{S}^*$ sous C1.

\textbf{\'Etape 2 : Approximation gaussienne locale.}
Par un d\'eveloppement de Taylor au second ordre de $\log P(\mathbf{D}_m
\mid \mathbf{S})$ autour de $\mathbf{S}^*$~:
\[
  \log P(\mathbf{D}_m \mid \mathbf{S}) \approx
  \log P(\mathbf{D}_m \mid \mathbf{S}^*)
  + \nabla \ell_m^\top (\mathbf{S} - \mathbf{S}^*)
  - \frac{1}{2}(\mathbf{S} - \mathbf{S}^*)^\top
  \bigl[-\nabla^2 \ell_m\bigr]
  (\mathbf{S} - \mathbf{S}^*)
\]
Par la LGN, $-\frac{1}{m}\nabla^2 \ell_m \xrightarrow{p.s.}
\mathcal{I}(\mathbf{S}^*)$ (C3, C4).

\textbf{\'Etape 3 : Gaussianisation du posterior.}
En substituant dans le posterior~:
\begin{align*}
  P(\mathbf{S} \mid \mathbf{D}_m)
  &\propto \exp\!\left(\log P(\mathbf{D}_m \mid \mathbf{S})\right)
  \cdot P(\mathbf{S}) \\
  &\approx \exp\!\left(-\frac{m}{2}(\mathbf{S}-\mathbf{S}^*)^\top
    \mathcal{I}(\mathbf{S}^*)
    (\mathbf{S}-\mathbf{S}^*)\right) \cdot P(\mathbf{S})
\end{align*}

\textbf{\'Etape 4 : Domination de l'a priori.}
Sous C2, $P(\mathbf{S}) > 0$ dans un voisinage de $\mathbf{S}^*$.
Quand $m \to \infty$, le terme exponentiel concentre la masse en
$\mathbf{S}^*$ \`a vitesse $m$, dominant l'a priori. Le posterior
converge vers~:
\[
  P(\mathbf{S} \mid \mathbf{D}_m)
  \xrightarrow{d} \mathcal{N}\!\left(\mathbf{S}^*,\,
  \frac{1}{m}\mathcal{I}(\mathbf{S}^*)^{-1}\right)
\]
Le r\'esultat suit par le changement de variable
$\sqrt{m}(\mathbf{S} - \mathbf{S}^*)$.
\end{proof}

\begin{corollaire}[Robustesse \`a l'a priori --- complet]
Soient $P_1(\mathbf{S})$ et $P_2(\mathbf{S})$ deux a priori RSI
satisfaisant C2. Alors~:
\[
  \mathrm{TV}\bigl(P_1(\mathbf{S} \mid \mathbf{D}_m),\,
  P_2(\mathbf{S} \mid \mathbf{D}_m)\bigr) \to 0
  \quad (m \to \infty)
\]
\begin{proof}
Les deux posteriors convergent vers
$\mathcal{N}(\mathbf{S}^*, m^{-1}\mathcal{I}(\mathbf{S}^*)^{-1})$
par le Th\'eor\`eme~2. Puisque la variation totale m\'etrise la
convergence faible sur $\mathbb{R}^d$, $\mathrm{TV}(P_i(\mathbf{S}
|\mathbf{D}_m), \mathcal{N}(\cdot)) \to 0$ pour $i = 1,2$, et le
r\'esultat suit par l'in\'egalit\'e triangulaire.
\end{proof}
\end{corollaire}

\subsection{Preuve du Th\'eor\`eme 3 : Convergence Monotone de l'ELBO}

\begin{theoreme}[ELBO Monotone --- complet]
Les mises \`a jour par ascension de coordonn\'ees RSI produisent une
s\'equence ELBO monotone non d\'ecroissante~:
$\mathrm{ELBO}_{t+1} \geq \mathrm{ELBO}_t$ pour tout $t \geq 0$.
\end{theoreme}

\begin{proof}
Rappelons la d\'ecomposition de l'ELBO~:
\[
  \mathrm{ELBO}(\phi) = \mathbb{E}_{Q_\phi}[\log P(\mathbf{D} \mid \mathbf{S})]
  - \mathrm{KL}[Q_\phi(\mathbf{S}) \| P(\mathbf{S})]
  = \log P(\mathbf{D}) - \mathrm{KL}[Q_\phi(\mathbf{S}) \| P(\mathbf{S} \mid \mathbf{D})]
\]
Puisque $\log P(\mathbf{D})$ est constant par rapport \`a $\phi$,
maximiser l'ELBO est \'equivalent \`a minimiser
$\mathrm{KL}[Q_\phi \| P(\cdot|\mathbf{D})]$.

\`A l'it\'eration $t$, la mise \`a jour pour le facteur $i$ r\'esout~:
\[
  \phi_i^{(t+1)} = \arg\min_{\phi_i}
  \mathrm{KL}\bigl[Q_{\phi_i}(s_i) \cdot \prod_{j \neq i}
  Q_{\phi_j^{(t)}}(s_j) \;\|\; P(\mathbf{S} \mid \mathbf{D})\bigr]
\]
La solution optimale est~:
\[
  Q_{\phi_i}^*(s_i) \propto \exp\!\left(
  \mathbb{E}_{Q_{-i}}\bigl[\log P(\mathbf{D}, \mathbf{S})\bigr]
  \right)
\]
Cette mise \`a jour ne peut que diminuer la divergence KL (ou la
laisser inchang\'ee), d'o\`u $\mathrm{ELBO}_{t+1} \geq \mathrm{ELBO}_t$.

Puisque l'ELBO est born\'e sup\'erieurement par $\log P(\mathbf{D})
< \infty$, la s\'equence converge. La limite est un point stationnaire
de l'ELBO.
\end{proof}

\begin{lemme}[Borne sur le gap variationnel]
Soit $Q_\phi^*$ l'optimum champ-moyen. Alors~:
\[
  \mathrm{KL}[Q_\phi^*(\mathbf{S}) \| P(\mathbf{S} \mid \mathbf{D})]
  \leq \sum_{i=1}^n \mathrm{KL}[Q_{\phi_i}^*(s_i) \| P(s_i \mid \mathbf{D})]
\]
\begin{proof}
Par la sous-additivit\'e de la divergence KL sous l'hypoth\`ese
d'ind\'ependance champ-moyen~:
$Q_\phi(\mathbf{S}) = \prod_i Q_{\phi_i}(s_i)$.
\end{proof}
\end{lemme}

%% ─────────────────────────────────────────────────────────────────
\section{D\'etails de Reproductibilit\'e}
\label{app:reproductibilite}

\subsection{Hyperparams de l'a priori}

Le Tableau~\ref{tab:hyperparams} reportent les hyperparams de l'a
priori $(\pi_i, \alpha_i, \beta_i, \sigma_i)$ utilis\'es pour chaque
r\`egle dans l'instanciation togolaise de RSI. Ces valeurs encodent
la connaissance institutionnelle sur les taux de conformit\'e
historiques et la stabilit\'e des r\`egles.

\begin{table}[H]
\centering
\caption{Hyperparams de l'a priori par r\`egle. $\pi_i$~: probabilit\'e
d'activation. $(\alpha_i, \beta_i)$~: a priori Beta sur la conformit\'e
($\mathbb{E}[c_i] = \alpha/(\alpha+\beta)$). $\sigma_i$~: \'ecart-type
du prior de d\'erive param\'etrique.}
\label{tab:hyperparams}
\renewcommand{\arraystretch}{1.3}
\begin{tabular}{llccccc}
\toprule
\textbf{R\`egle} & \textbf{Description} & $\pi_i$ & $\alpha_i$ & $\beta_i$
  & $\mathbb{E}[c_i]$ & $\sigma_i$ \\
\midrule
R1\_TVA & TVA 18\%, seuil 60M / 100M FCFA & 0,92 & 8,0 & 2,0 & 0,80 & 0,05 \\
R2\_IS  & IS 29\%, seuil 100M FCFA         & 0,88 & 6,0 & 4,0 & 0,60 & 0,03 \\
R3\_IMF & IMF 1\% du CA                     & 0,85 & 9,0 & 1,5 & 0,86 & 0,02 \\
R4\_TPU & Forfait secteur informel          & 0,70 & 3,0 & 7,0 & 0,30 & 0,15 \\
\bottomrule
\end{tabular}
\end{table}

\noindent\textbf{Justification des choix de l'a priori.}
\begin{itemize}[leftmargin=2em]
  \item \textbf{R1\_TVA}~: Conformit\'e a priori \'elev\'ee
        ($\mathbb{E}[c_i]=0{,}80$) car les entreprises assujettis \`a la
        TVA au Togo d\'eposent g\'en\'eralement leurs d\'eclarations,
        m\^eme si le montant est sous-estim\'e. Faible d\'erive
        ($\sigma=0{,}05$) car le taux de 18\% est stable sur la p\'eriode.
  \item \textbf{R2\_IS}~: Conformit\'e mod\'er\'ee ($0{,}60$) refl\'etant
        la complexit\'e du calcul IS et les taux d'\'evasion plus \'elev\'es
        parmi les grandes entreprises.
  \item \textbf{R3\_IMF}~: Conformit\'e \'elev\'ee ($0{,}86$) car l'IMF
        est plus simple \`a calculer (1\% du CA) et plus difficile \`a
        \'eviter.
  \item \textbf{R4\_TPU}~: Faible conformit\'e ($0{,}30$) refl\'etant
        l'informalit\'e bien document\'ee du secteur des micro-entreprises
        en Afrique subsaharienne. D\'erive \'elev\'ee ($\sigma=0{,}15$)
        refl\'etant la forte incertitude sur ce segment.
\end{itemize}

\subsection{Param\`etres de l'inf\'erence variationnelle}

\begin{table}[H]
\centering
\caption{Hyperparams de l'inf\'erence variationnelle.}
\renewcommand{\arraystretch}{1.2}
\begin{tabular}{lcc}
\toprule
\textbf{Param\`etre} & \textbf{Valeur} & \textbf{Justification} \\
\midrule
Taux d'apprentissage $\eta$ & 1,5 & Amplifie la mise \`a jour post\'erieure \\
It\'erations max & 150 & Convergence empirique $\leq$ 10 it\'erations \\
Tol\'erance de convergence $\varepsilon$ & $10^{-5}$ & Seuil de variation ELBO \\
Seuil de d\'ecision $\tau$ & 0,551 & Optimis\'e pour le F1 sur l'ensemble de validation \\
\bottomrule
\end{tabular}
\end{table}

\subsection{Configurations des baselines}

\begin{table}[H]
\centering
\caption{Configurations des mod\`eles baseline.}
\renewcommand{\arraystretch}{1.2}
\begin{tabular}{lll}
\toprule
\textbf{Mod\`ele} & \textbf{Param\`etre} & \textbf{Valeur} \\
\midrule
XGBoost & n\_estimators & 150 \\
        & max\_depth & 4 \\
        & learning\_rate & 0,1 \\
        & subsample & 0,8 \\
        & random\_state & 42 \\[4pt]
MLP     & hidden\_layer\_sizes & (64, 32, 16) \\
        & activation & ReLU \\
        & max\_iter & 200 \\
        & early\_stopping & Vrai \\
        & validation\_fraction & 0,1 \\
        & random\_state & 42 \\[4pt]
SBR     & Seuil TVA (2022--2024) & 60M FCFA \\
        & Seuil TVA (2025) & 100M FCFA \\
        & Seuil d'alerte retard & 90 jours \\
\bottomrule
\end{tabular}
\end{table}

\subsection{Environnement de calcul}

Les exp\'eriences ont \'et\'e ex\'ecut\'ees sur un ordinateur de
bureau standard (Windows 11, processeur Intel Core, 16 Go de RAM).
Tout le code est impl\'ement\'e en Python 3.11 en utilisant NumPy,
SciPy et scikit-learn. Aucun GPU n'a \'et\'e utilis\'e. La dur\'ee
du pipeline exp\'erimental complet est d'environ 5 \`a 10 minutes.

\subsection{Graines al\'eatoires}

Toutes les exp\'eriences utilisent \texttt{numpy.random.seed(42)}
pour la reproductibilit\'e. La g\'en\'eration du dataset utilise
\texttt{numpy.random.default\_rng(42)}.

%% ─────────────────────────────────────────────────────────────────
\begin{thebibliography}{9}

\bibitem[Gabaix(2009)]{gabaix2009power}
Gabaix, X. (2009).
Power laws in economics and finance.
\textit{Annual Review of Economics}, 1(1):255--294.

\bibitem[Blei et al.(2017)]{blei2017variational}
Blei, D. M., Kucukelbir, A., et McAuliffe, J. D. (2017).
Variational inference: A review for statisticians.
\textit{Journal of the American Statistical Association}, 112(518):859--877.

\bibitem[van der Vaart(2000)]{van2000asymptotic}
van der Vaart, A. W. (2000).
\textit{Asymptotic Statistics}.
Cambridge University Press.

\end{thebibliography}

\end{document}